\subsection{A genuine common-skeleton pullback need not admit forward completion}
\label{arch136:coskeleton-direction}

\paragraph{ARCH-136-COSKELETON-FORWARD.}
Let \(C\) be an admissible strict differential PD theory over
\(S=\Gamma_{\mathbb Z}(T)\), with \(C_0=S\), and let \(r\geq0\)
be finite. Write
\[
 U_r=\operatorname{cosk}_r\operatorname{sk}_r C,
 \qquad P_r=C\times_{U_r}C,
 \qquad \pi_0,\pi_1:P_r\longrightarrow C.
 \tag{CF.1}
\]
There is no admissible theory \(E\) and no strict diagram
\[
 P_r\xrightarrow{h}E
       \mathrel{\substack{\xrightarrow{g_0}\\[-.3ex]
                           \xrightarrow[g_1]{}}}C
 \quad\text{with}\quad g_i h=\pi_i\ (i=0,1).
 \tag{CF.2}
\]
Thus a genuine fiber product of two admissible theories over a
common finite skeleton can fail the prescribed-map forward
completion problem, even when the two theories are identical.
This does not obstruct an admissible source mapping into the
pullback: the diagonal \(C\to P_r\) already is such a map.

\paragraph{Proof.}
Put \(s=r+1\) and use the simplicial sphere
\[
                    X=\Delta[s]/\partial\Delta[s].
 \tag{CF.3}
\]
Its \(r\)-skeleton is the one-point simplicial set. By the
skeleton--coskeleton adjunction, applied to every bigraded
component and its compatible algebra operations,
\[
 U_r(X)=C(\operatorname{sk}_r X)=C_0=S.
 \tag{CF.4}
\]
Consequently restriction to \(U_r(X)\) kills every cochain of
positive form degree, and sections preserve the pullback:
\[
                     P_r(X)=C(X)\times_S C(X).
 \tag{CF.5}
\]

Suppose (CF.2) exists. By AQ-136-COMPARISON, choose one weight
\(w\) sufficiently large that the natural coefficient comparisons
in degree \(s\) hold for both \(C(X)\) and \(E(X)\). These
comparisons identify
\[
 H^{w,s}(C(X))\simeq
 H^s(X;\mathbb Z)T^{[w]}\simeq\mathbb ZT^{[w]},
 \qquad
 H^{w,s}(E(X))\simeq\mathbb ZT^{[w]}.
 \tag{CF.6}
\]
Both \(g_0\) and \(g_1\) preserve the specified coefficient
map. Naturality of this same comparison therefore gives
\[
          H^{w,s}(g_0(X))=H^{w,s}(g_1(X)).
 \tag{CF.7}
\]
No choice of cochain representative enters this equality.

Choose a cocycle \(z\in C^{w,s}(X)\) representing a nonzero
class in (CF.6). Because \(s>0\), (CF.4) makes
\((0,z)\) a cocycle of \(P_r^{w,s}(X)\). Its images under
\(\pi_0,\pi_1\) are respectively zero and the nonzero class
\([z]\). But applying (CF.7) to \(H^{w,s}(h(X))[(0,z)]\)
would make those images equal. This is a contradiction, proving
the statement. The diagonal source map in the final sentence
follows directly from the pullback universal property.
\(\square\)

\paragraph{Exact consequence for the degree-two proposal.}
For \(r=1\), take \(C=B\), the standard admissible integral
theory. This is precisely the constrained pullback
\(B\times_{\operatorname{cosk}_1\operatorname{sk}_1 B}B\).
Even if its high-weight simplexwise vertical calculation is
\(H^0=S\), \(H^1=0\), and free \(H^2\) identified with the
first cohomology of the simplex's edge graph, no successive
under-pullback attachments can make it admissible while extending
both projections. The obstruction is already the different
primary degree-two images on the simplicial two-sphere, before
one asks whether the coefficient tail has an \(S\)-linear cycle
section or a nonfree Koszul attachment.

The higher-skeleton prescribed-map example gives another check:
with its parameter \(d\geq r+1\), its two maps to \(B\) agree
through simplicial degree \(r\), and hence factor jointly through
this pullback. A forward completion would also contradict that
example. The direct sphere proof above avoids importing it and
shows that the failure holds for every admissible \(C\) with
\(C_0=S\). Neither proof says that the original objects lack a
common source; the identical-target case has the diagonal source
explicitly.
